%% ========================================================================
%% CAPÍTULO 3 - RESULTADOS E DISCUSSÃO
%% ========================================================================

\capitulo{resultados}{Resultados e Discussão}

Este capítulo apresenta e discute os resultados obtidos através da análise dos dados coletados. Os resultados são organizados em seções que abordam a caracterização dos padrões de mobilidade, avaliação dos sistemas de transporte e proposição de estratégias para mobilidade sustentável.

\section{Caracterização dos Padrões de Mobilidade}

\subsection{Perfil Socioeconômico dos Respondentes}

A amostra de 450 respondentes apresentou distribuição equilibrada em relação às variáveis de estratificação. A \refTab{tab:perfil-respondentes} apresenta as características socioeconômicas dos participantes.

\begin{table}[htb]
    \centering
    \caption{Perfil socioeconômico dos respondentes}
    \label{tab:perfil-respondentes}
    \begin{tabular}{llrr}
        \toprule
        \textbf{Variável} & \textbf{Categoria} & \textbf{n} & \textbf{\%} \\
        \midrule
        \multirow{2}{*}{Localização} & DF & 315 & 70,0 \\
                                     & Entorno & 135 & 30,0 \\
        \midrule
        \multirow{3}{*}{Renda Familiar} & Até 3 SM & 180 & 40,0 \\
                                       & 3-6 SM & 158 & 35,1 \\
                                       & Acima de 6 SM & 112 & 24,9 \\
        \midrule
        \multirow{3}{*}{Faixa Etária} & 18-30 anos & 158 & 35,1 \\
                                      & 31-50 anos & 202 & 44,9 \\
                                      & Acima de 50 anos & 90 & 20,0 \\
        \midrule
        \multirow{2}{*}{Sexo} & Feminino & 248 & 55,1 \\
                              & Masculino & 202 & 44,9 \\
        \bottomrule
    \end{tabular}
\end{table}

\subsection{Distribuição Modal dos Deslocamentos}

A análise dos padrões de deslocamento revelou a predominância do transporte individual motorizado, conforme apresentado na \refFig{fig:distribuicao-modal}.

\begin{figure}[htb]
    \centering
    \includegraphics[width=0.8\textwidth]{img/distribuicao-modal.png}
    \caption{Distribuição modal dos deslocamentos na Região Metropolitana de Brasília}
    \label{fig:distribuicao-modal}
\end{figure}

Os resultados indicam que:
\begin{itemize}
    \item \textbf{Automóvel particular}: 65,2\% dos deslocamentos;
    \item \textbf{Transporte público}: 24,8\% dos deslocamentos;
    \item \textbf{Caminhada}: 7,3\% dos deslocamentos;
    \item \textbf{Bicicleta}: 2,1\% dos deslocamentos;
    \item \textbf{Outros}: 0,6\% dos deslocamentos.
\end{itemize}

Esta distribuição confirma a hipótese de predominância do transporte individual motorizado e está alinhada com estudos anteriores realizados na região \cite{silva2019, oliveira2021}.

\subsection{Características dos Deslocamentos}

\subsubsection{Motivos de Viagem}

A \refTab{tab:motivos-viagem} apresenta a distribuição dos deslocamentos por motivo, evidenciando a importância das viagens trabalho-residência.

\begin{table}[htb]
    \centering
    \caption{Distribuição dos deslocamentos por motivo de viagem}
    \label{tab:motivos-viagem}
    \begin{tabular}{lrr}
        \toprule
        \textbf{Motivo} & \textbf{Frequência} & \textbf{Percentual} \\
        \midrule
        Trabalho & 1.890 & 42,0 \\
        Estudo & 810 & 18,0 \\
        Compras/Serviços & 720 & 16,0 \\
        Lazer & 450 & 10,0 \\
        Saúde & 315 & 7,0 \\
        Outros & 315 & 7,0 \\
        \midrule
        \textbf{Total} & \textbf{4.500} & \textbf{100,0} \\
        \bottomrule
    \end{tabular}
\end{table}

\subsubsection{Tempos de Deslocamento}

Os tempos médios de deslocamento variam significativamente entre os modais, conforme apresentado na \refFig{fig:tempos-deslocamento}.

\begin{figure}[htb]
    \centering
    \includegraphics[width=0.9\textwidth]{img/tempos-deslocamento.png}
    \caption{Tempos médios de deslocamento por modal de transporte}
    \label{fig:tempos-deslocamento}
\end{figure}

Os resultados mostram que:
\begin{itemize}
    \item Usuários de transporte público gastam, em média, 78 minutos por deslocamento;
    \item Usuários de automóvel gastam, em média, 35 minutos por deslocamento;
    \item A diferença de tempo é um fator determinante na escolha modal.
\end{itemize}

\section{Avaliação dos Sistemas de Transporte}

\subsection{Qualidade do Transporte Público}

A avaliação da qualidade do transporte público foi realizada através de uma escala Likert de 5 pontos. A \refTab{tab:avaliacao-tp} apresenta as médias obtidas para cada dimensão avaliada.

\begin{table}[htb]
    \centering
    \caption{Avaliação da qualidade do transporte público}
    \label{tab:avaliacao-tp}
    \begin{tabular}{lrrr}
        \toprule
        \textbf{Dimensão} & \textbf{Média} & \textbf{Desvio Padrão} & \textbf{Classificação} \\
        \midrule
        Pontualidade & 2,1 & 0,8 & Ruim \\
        Frequência & 2,3 & 0,9 & Ruim \\
        Conforto & 2,0 & 0,7 & Ruim \\
        Segurança & 2,4 & 1,0 & Ruim \\
        Limpeza & 2,8 & 0,9 & Regular \\
        Acessibilidade & 2,2 & 0,8 & Ruim \\
        Integração & 1,9 & 0,7 & Ruim \\
        \midrule
        \textbf{Média Geral} & \textbf{2,2} & \textbf{0,5} & \textbf{Ruim} \\
        \bottomrule
    \end{tabular}
\end{table}

Os resultados evidenciam deficiências significativas em todas as dimensões avaliadas, com destaque para os problemas de integração entre modais e conforto dos veículos.

\subsection{Infraestrutura para Modos Não Motorizados}

A avaliação da infraestrutura para pedestres e ciclistas revelou inadequações importantes:

\begin{itemize}
    \item \textbf{Calçadas}: 68\% dos respondentes consideram inadequadas;
    \item \textbf{Ciclovias}: Apenas 12\% da malha viária possui infraestrutura cicloviária;
    \item \textbf{Segurança}: 74\% consideram inseguro caminhar ou pedalar;
    \item \textbf{Conectividade}: Falta de continuidade nas rotas.
\end{itemize}

\subsection{Análise de Correlação}

A análise de correlação entre variáveis revelou associações significativas (p < 0,05):

\begin{itemize}
    \item Renda familiar e uso do automóvel (r = 0,68);
    \item Tempo de deslocamento e satisfação com transporte público (r = -0,52);
    \item Idade e uso de modos não motorizados (r = -0,34);
    \item Distância residência-trabalho e escolha modal (r = 0,45).
\end{itemize}

\section{Impactos Ambientais e Socioeconômicos}

\subsection{Emissões de Gases Poluentes}

O cálculo das emissões de CO₂ por modal foi realizado utilizando fatores de emissão do \sigla{IPCC}. A \refFig{fig:emissoes-co2} apresenta os resultados obtidos.

\begin{figure}[htb]
    \centering
    \includegraphics[width=0.8\textwidth]{img/emissoes-co2.png}
    \caption{Emissões de CO₂ por passageiro-km por modal de transporte}
    \label{fig:emissoes-co2}
\end{figure}

Os resultados mostram que:
\begin{itemize}
    \item Automóvel particular: 120 g CO₂/pass-km;
    \item Ônibus: 45 g CO₂/pass-km;
    \item Metrô: 25 g CO₂/pass-km;
    \item Bicicleta: 0 g CO₂/pass-km.
\end{itemize}

\subsection{Custos Socioeconômicos}

A análise dos custos socioeconômicos considerou:

\begin{enumerate}
    \item \textbf{Custos diretos}: Combustível, manutenção, tarifas;
    \item \textbf{Custos indiretos}: Tempo perdido, acidentes, poluição;
    \item \textbf{Custos de infraestrutura}: Construção e manutenção viária.
\end{enumerate}

A \refTab{tab:custos-socioeconomicos} apresenta os custos totais por modal.

\begin{table}[htb]
    \centering
    \caption{Custos socioeconômicos por modal (R\$/pass-km)}
    \label{tab:custos-socioeconomicos}
    \begin{tabular}{lrrrr}
        \toprule
        \textbf{Modal} & \textbf{Diretos} & \textbf{Indiretos} & \textbf{Infraestrutura} & \textbf{Total} \\
        \midrule
        Automóvel & 0,45 & 0,32 & 0,18 & 0,95 \\
        Ônibus & 0,25 & 0,12 & 0,08 & 0,45 \\
        Metrô & 0,35 & 0,05 & 0,15 & 0,55 \\
        Bicicleta & 0,02 & 0,01 & 0,03 & 0,06 \\
        Caminhada & 0,00 & 0,01 & 0,02 & 0,03 \\
        \bottomrule
    \end{tabular}
\end{table}

\section{Proposição de Estratégias}

Com base nos resultados obtidos, foram propostas estratégias integradas organizadas em cinco eixos principais:

\subsection{Eixo 1: Melhoria do Transporte Público}

\subsubsection{Expansão da Rede}
\begin{itemize}
    \item Extensão do metrô para cidades do entorno;
    \item Implementação de \sigla{BRT} nos principais corredores;
    \item Criação de linhas alimentadoras eficientes.
\end{itemize}

\subsubsection{Melhoria Operacional}
\begin{itemize}
    \item Aumento da frequência nos horários de pico;
    \item Implementação de sistema de informação ao usuário;
    \item Modernização da frota com veículos acessíveis.
\end{itemize}

\subsection{Eixo 2: Integração Modal}

A \refFig{fig:integracao-modal} apresenta o conceito de integração proposto.

\begin{figure}[htb]
    \centering
    \includegraphics[width=0.9\textwidth]{img/integracao-modal.png}
    \caption{Modelo conceitual de integração modal}
    \label{fig:integracao-modal}
\end{figure}

\subsubsection{Integração Física}
\begin{itemize}
    \item Construção de terminais intermodais;
    \item Conexões diretas entre estações;
    \item Facilidades para transferência.
\end{itemize}

\subsubsection{Integração Tarifária}
\begin{itemize}
    \item Bilhete único para todos os modais;
    \item Tarifas integradas com desconto;
    \item Sistema de pagamento eletrônico.
\end{itemize}

\subsection{Eixo 3: Infraestrutura para Modos Não Motorizados}

\subsubsection{Rede Cicloviária}
A proposta prevê a criação de uma rede cicloviária integrada, conforme equação de densidade ótima:

\begin{equation}
D_{ciclo} = \frac{L_{ciclo}}{A_{urbana}} \geq 0,5 \text{ km/km}^2
\label{eq:densidade-ciclo}
\end{equation}

Onde $D_{ciclo}$ é a densidade cicloviária, $L_{ciclo}$ é o comprimento total de ciclovias e $A_{urbana}$ é a área urbana.

\subsubsection{Infraestrutura para Pedestres}
\begin{itemize}
    \item Melhoria das calçadas existentes;
    \item Criação de travessias seguras;
    \item Iluminação adequada;
    \item Acessibilidade universal.
\end{itemize}

\subsection{Eixo 4: Gestão da Demanda}

\subsubsection{Políticas de Desestímulo ao Automóvel}
\begin{itemize}
    \item Cobrança de estacionamento em áreas centrais;
    \item Pedágio urbano em horários de pico;
    \item Restrição de circulação por rodízio.
\end{itemize}

\subsubsection{Incentivos aos Modos Sustentáveis}
\begin{itemize}
    \item Subsídios para transporte público;
    \item Programas de bike-sharing;
    \item Campanhas de conscientização.
\end{itemize}

\subsection{Eixo 5: Governança e Financiamento}

\subsubsection{Arranjo Institucional}
\begin{itemize}
    \item Criação de autoridade metropolitana de transporte;
    \item Coordenação entre diferentes níveis de governo;
    \item Participação social nas decisões.
\end{itemize}

\subsubsection{Fontes de Financiamento}
\begin{itemize}
    \item Recursos federais (PAC Mobilidade);
    \item Parcerias público-privadas;
    \item Contribuição de melhoria;
    \item Fundos internacionais de desenvolvimento sustentável.
\end{itemize}

\section{Indicadores de Monitoramento}

Para acompanhar a implementação das estratégias propostas, foram definidos indicadores de monitoramento organizados em três dimensões:

\subsection{Indicadores de Resultado}
\begin{itemize}
    \item Participação do transporte público na matriz modal;
    \item Tempo médio de deslocamento;
    \item Emissões de CO₂ per capita;
    \item Índice de satisfação dos usuários.
\end{itemize}

\subsection{Indicadores de Processo}
\begin{itemize}
    \item Quilômetros de ciclovias construídas;
    \item Número de terminais intermodais;
    \item Frequência média do transporte público;
    \item Cobertura espacial dos serviços.
\end{itemize}

\subsection{Indicadores de Impacto}
\begin{itemize}
    \item Redução de acidentes de trânsito;
    \item Melhoria da qualidade do ar;
    \item Inclusão social e acessibilidade;
    \item Desenvolvimento econômico local.
\end{itemize}

\section{Discussão dos Resultados}

Os resultados obtidos confirmam as hipóteses iniciais da pesquisa e estão alinhados com a literatura internacional sobre mobilidade urbana sustentável. A predominância do transporte individual motorizado (65,2\%) é consistente com outras regiões metropolitanas brasileiras, mas superior à média de cidades com sistemas de transporte mais desenvolvidos.

A baixa avaliação do transporte público (média 2,2 em escala de 5 pontos) evidencia a necessidade urgente de investimentos em qualidade e expansão dos serviços. Os problemas identificados - baixa frequência, falta de pontualidade e deficiências na integração - são obstáculos significativos para a migração modal.

As estratégias propostas seguem as melhores práticas internacionais, priorizando a hierarquia de mobilidade urbana e a integração entre diferentes modais. A implementação bem-sucedida dessas estratégias requer coordenação institucional, recursos financeiros adequados e participação social.

Os indicadores de monitoramento propostos permitirão avaliar o progresso das políticas implementadas e realizar ajustes necessários. A experiência internacional mostra que a transição para sistemas de mobilidade mais sustentáveis é um processo gradual que requer persistência e adaptação contínua.